\documentclass{article}
\usepackage{graphicx} % Required for inserting images
\usepackage[margin=1in]{geometry}
\usepackage{enumitem}
\usepackage{amsmath}
\usepackage{hyperref}

\title{Personal Finance Manager: A Desktop-Based Financial\\ Tracking and Budgeting System \\ \large Project Proposal for UCS503P}
\author{
  Arjun Puri\thanks{Roll No: 1024030022, Email: apuri2\_be24@thapar.edu} \and
  Anshdeep Singh \thanks{Roll No: 1024030105}
}
\date{August 2026 \\ \small Submitted to: Ms. Parmveer  \\ Thapar Institute of Engineering and Technology}

\begin{document}
\maketitle

\section{Higher-Order Goal --- secondary framing}
This project aims to improve personal financial literacy and money discipline by reducing the effort required to track income, control spending, and reach savings targets. While the underlying idea applies to any household budgeting context, the immediate engineering objective is to translate everyday financial record-keeping into a measurable, deployable software solution -- one that a student or family can trust with real financial data.

\section{Time-to-Value --- primary heuristic}
\begin{itemize}[noitemsep]
  \item We deliver meaningful increments quickly by building the transaction-and-database core first and validating it before any reporting or chart features are attempted.
  \item We prioritize fast feedback over perfection: a working login $\rightarrow$ add-transaction $\rightarrow$ updated-balance loop ships early, then budgets, goals and recurring items layer on top of it.
  \item Success is defined by what can be delivered and measured in the first development phase, then improved based on self-testing and peer feedback.
\end{itemize}

\section{Problem Statement}
Students and individuals typically receive income from multiple sources (allowance, stipends, part-time work) and spend across scattered categories -- food, travel, subscriptions, shopping, education and bills. Without a structured system, it is difficult to know where money is going or whether monthly spending is under control. Common pain points include:
\begin{itemize}[noitemsep]
  \item \textbf{Fragmentation:} transactions are scattered across bank SMS alerts, UPI apps and cash, with no single source of truth.
  \item \textbf{Low visibility:} no simple, immediate view of how actual spending compares against a monthly budget.
  \item \textbf{No goal tracking:} savings targets (e.g. for a laptop or a trip) are rarely tracked with a visible progress measure.
  \item \textbf{Manual, repetitive entry:} recurring costs such as rent, subscriptions and bills are re-entered from memory each month, which is error-prone.
\end{itemize}
\textit{Why this matters now (the need):} for a student managing a limited and irregular monthly budget, even small improvements in visibility and habit-tracking noticeably improve financial discipline and reduce overspending -- and unlike a generic notes app, a purpose-built tool can enforce consistency and catch mistakes automatically.

\section{Proposed Solution --- Software Proposition}
\subsection{Overview}
A desktop ``Finance Manager'' application with:
\begin{itemize}[noitemsep]
  \item \textbf{Transaction entry:} log income/expenses with amount, category, date, description.
  \item \textbf{Budget engine:} monthly or category-wise spending limits with real-time remaining balance.
  \item \textbf{Savings goals:} target amount, current saved amount, deadline and visible progress.
  \item \textbf{Recurring expenses:} rent, subscriptions and bills scheduled by frequency, surfaced automatically.
  \item \textbf{Dashboard and reports:} balance, income/expense summaries, category and monthly charts.
\end{itemize}

\subsection{Core Workflow}
\begin{enumerate}[noitemsep]
  \item User registers/logs in.
  \item Dashboard loads balance, income, expenses, budget status and goal progress.
  \item User adds a transaction (income/expense, amount, category, date, description).
  \item System validates the input (amount, date, required fields).
  \item Transaction is written to MySQL.
  \item Balance, category totals, budget usage and savings are recalculated and the dashboard refreshes.
\end{enumerate}

\subsection{Operational Constraints for an Educational Setting}
\begin{itemize}[noitemsep]
  \item Keep the solution lightweight enough to run on a standard student laptop with a local MySQL instance.
  \item Avoid dependence on advanced research algorithms; focus on correctness of financial calculations.
  \item Design the data layer so it can later be exposed behind a web API without redesigning the schema.
\end{itemize}

\section{Solution Approach --- Engineering Focus}
This is a software engineering course project, so the emphasis is on correct, modular application design rather than algorithmic novelty.

\subsection{Financial Calculations --- non-research}
All dashboard figures are deterministic, auditable formulas -- important for a finance tool, since a user must be able to \emph{trust} every number shown:
\begin{itemize}[noitemsep]
  \item Balance = Total Income $-$ Total Expenses
  \item Savings Rate = (Savings / Income) $\times$ 100
  \item Budget Remaining = Budget Limit $-$ Actual Spending
  \item Goal Progress = (Current Saved / Target Amount) $\times$ 100
\end{itemize}

\subsection{Application Architecture --- 3-tier}
\begin{itemize}[noitemsep]
  \item \textbf{Presentation:} Tkinter GUI -- login, dashboard, transaction forms, budget/goal views.
  \item \textbf{Business logic:} Python modules for validation, calculations, recurring-expense scheduling.
  \item \textbf{Data:} MySQL via PyMySQL -- users, categories, transactions, budgets, savings goals, recurring expenses.
\end{itemize}

\subsection{Testing and Delivery Discipline}
\begin{itemize}[noitemsep]
  \item Unit tests for calculation functions (balance, budget remaining, goal progress) run before every merge.
  \item Manual regression checklist for the transaction $\rightarrow$ dashboard $\rightarrow$ report path before each phase is marked complete.
  \item A lightweight script-based build/run check stands in for full CI/CD now, with a clear upgrade path to a real CI runner if the project moves to a web/Flask version.
\end{itemize}

\section{Evaluation Criterion --- measurable and attributable}
We define success using measurable metrics tied to the core workflow, not subjective polish -- this is what makes the result convincing rather than anecdotal.

\subsection{Primary Evaluation Metric}
\textbf{Dashboard Refresh Latency (DRL):} time between saving a transaction and the dashboard/report views reflecting the updated balance, budget and charts.
\begin{itemize}[noitemsep]
  \item Target: DRL $\leq$ 1 second on a standard student laptop for up to 5{,}000 transactions.
  \item Attribution: measured by timestamping the save call and the completed UI refresh.
\end{itemize}

\subsection{Secondary Evaluation Metrics}
\begin{itemize}[noitemsep]
  \item \textbf{Calculation accuracy:} 100\% agreement between system-computed figures and hand-calculated values on a fixed test dataset.
  \item \textbf{Task completion time:} average time for a first-time user to add a transaction and locate their updated balance.
  \item \textbf{Reliability:} percentage of test sessions completed without a crash or unhandled exception.
  \item \textbf{Report usefulness:} peer-reviewer rating (1--5) of whether the report clearly answers ``where did my money go?''
\end{itemize}

\subsection{Pilot Validation Plan}
\begin{itemize}[noitemsep]
  \item Recruit 3--5 peer testers to use the application for a simulated month of transactions.
  \item Compare DRL and calculation accuracy against the fixed test dataset across development phases.
  \item Collect qualitative feedback on dashboard clarity and fold it into the UI-polishing phase.
\end{itemize}

\section{Scalability --- with foresight}
\begin{enumerate}[noitemsep]
  \item \textbf{Scalable (theoretically):}
  \begin{itemize}[noitemsep]
    \item Decoupled presentation / business-logic / data layers, so any layer can be replaced independently (e.g. Tkinter $\rightarrow$ Flask/React).
    \item Indexing frequently-queried columns (\texttt{user\_id}, \texttt{date}, \texttt{category\_id}) so queries stay fast as transaction volume grows.
    \item Schema designed to support multiple accounts/wallets per user (cash, bank, UPI) without breaking existing calculations.
  \end{itemize}
  \item \textbf{Operationally deployable (practically):}
  \begin{itemize}[noitemsep]
    \item Packaged as a standalone executable for easy distribution and demo.
    \item Business logic already lives outside the UI, so the same schema and calculations move behind a Flask REST API with minimal rework.
    \item A hosted managed database is a drop-in replacement for the local database when moving to a multi-user version.
  \end{itemize}
\end{enumerate}
\[
\text{Desktop MVP} \rightarrow \text{Web App (Flask)} \rightarrow \text{Multi-user Cloud} \rightarrow \text{Mobile App} \rightarrow \text{Family/Shared Budgets}
\]

\section{Engine Availability Heuristic}
A mature, well-documented set of ``engines'' is available, so effort goes into data modelling and correctness rather than infrastructure:
\begin{itemize}[noitemsep]
  \item \textbf{Python} -- application logic, validation, orchestration.
  \item \textbf{Tkinter} -- desktop GUI for all screens.
  \item \textbf{MySQL + PyMySQL} -- persistent relational storage and connectivity.
  \item \textbf{Matplotlib} -- monthly trend and category-distribution charts.
  \item \textbf{datetime / csv} -- recurring-expense scheduling and optional export.
\end{itemize}

\section{Project Scope and Deliverables --- iteration-friendly}
\subsection{Initial Deliverable --- first iteration}
\begin{itemize}[noitemsep]
  \item Database schema: users, categories, transactions tables
  \item Login/registration and session handling
  \item Add/view transaction workflow with validation
  \item Dashboard: balance, income, expenses (auto-recalculated)
  \item Unit tests for the core calculation functions
\end{itemize}

\subsection{Subsequent Deliverables}
\begin{itemize}[noitemsep]
  \item Category management and budget engine
  \item Savings goals with progress tracking
  \item Recurring expenses with due-date surfacing
  \item Reports \& charts (Matplotlib) and CSV export
  \item UI polishing, error handling, full regression pass
\end{itemize}

\section{Risks and Mitigations}
\begin{itemize}[noitemsep]
  \item \textbf{Financial data sensitivity} -- hash passwords, minimize stored fields, keep the database local/self-hosted for the pilot (builds user trust).
  \item \textbf{Calculation errors compounding} -- unit-test core formulas before building dashboard/report features on top of them.
  \item \textbf{Scope creep} -- follow the phased plan strictly; optional features (export, notifications) last.
  \item \textbf{Cross-user data leakage} -- enforce \texttt{user\_id} scoping on every query; test explicitly for it.
\end{itemize}

\section{Summary}
Finance Manager is a modular, database-driven personal finance application that lets a user track income and expenses, manage budgets and savings goals, handle recurring payments, and understand their financial behaviour through reports and a visual dashboard. Its strong points are what make it convincing as a course project: transparent, auditable calculations instead of black-box logic; a layered architecture that keeps the system scalable from a single-user desktop MVP to a multi-user web application without a redesign; and evaluation criteria (dashboard latency, calculation accuracy) that are measurable rather than anecdotal. Together these make the system not just functional, but trustworthy and built with clear headroom to grow.

\end{document}